\documentclass[10pt,a4paper]{amsart}
\usepackage[margin=19mm]{geometry}
\usepackage{amsmath,amssymb,mathtools}
\usepackage[scr=rsfs,cal=euler]{mathalfa}
\usepackage{xfrac}
\usepackage[unicode,hidelinks,bookmarks=false]{hyperref}
\newcommand{\ie}{i.e.\ }
\newcommand{\cf}{cf.\ }
\newcommand{\resp}{respectively\ }
\newcommand{\Subsection}{\subsection}
\providecommand{\nobreakdash}{\nobreak\hbox{-}}
\setlength{\emergencystretch}{2em}
\multlinegap0pt
\usepackage{bm}
\usepackage{bbm}
\usepackage{dsfont}
\usepackage{xspace}
\usepackage{cancel}
\usepackage{mathtools}
\usepackage[shortlabels]{enumitem}
\usepackage{amsthm}
\makeatletter
\@ifundefined{prop}{\newtheorem{prop}{Proposition}}{}
\makeatother
\DeclareMathOperator{\Cc}{\mathbb{C}}
\DeclareMathOperator{\inv}{\mathrm{inv}}
\renewcommand{\epsilon}{\varepsilon}
\title[$n$-Valued Groups, Kronecker Sums, and Wendt's Matrices]{$n$-Valued Groups, Kronecker Sums, and Wendt's Matrices}
\author{Victor Buchstaber, Mikhail Kornev}
\date{Source: arXiv:2505.04296v2 (CC BY 4.0)}
\begin{document}
\maketitle

\noindent\emph{Source and licence.} $n$-Valued Groups, Kronecker Sums, and Wendt's Matrices. Version arXiv:2505.04296v2 (CC BY 4.0), distributed under the Creative Commons Attribution 4.0 International licence, as recorded in the arXiv licence metadata for this exact version and shown on the abstract page https://arxiv.org/abs/2505.04296v2. Full rights notice: licenses/arxiv-2505.04296-CC-BY-4.0.txt.\par\smallskip

\noindent\textit{Editorial scope.} Counted unit: the Prop whose source label is \texttt{prop10} in the article (source0.tex lines 874--887, source label \texttt{prop10}), delivered together with the article's own complete original proof of it (source0.tex lines 889--913, one \texttt{proof} environment of 25 lines). No mathematics is written, repaired or re-derived: statement and proof are the article's own text line for line. Required context retained uncounted: coroll1 (lines 789--793); p\_n (lines 869--871). Every definition, display and label of those lines is retained unchanged. Omitted from this package and not redistributed: 1--788, 794--868, 872--873, 888--888, 914--1752. Nothing inside a retained excerpt is omitted or altered: every retained body is delivered exactly as the article's lines are written, so no omission receipt of any kind accompanies this package. Citations: the retained excerpts carry no citation command at all, so no bibliography entry of the article is transcribed and no reference is redistributed. Reference adaptation: none was needed and none was made; every reference inside the delivered unit resolves inside it, so no unresolved reference or citation is produced. Printed slips kept literal: no printed word, symbol, hypothesis or number of the counted statement or of its proof was altered. Layout and class: the article's own class is \texttt{extarticle} with packages including geometry, fontenc, newtxmath, ebgaramond, ebgaramond-maths, euscript, xcolor, graphicx, caption, microtype, amsthm, hyperref, breakurl, enumerate; the wrapper is the lane's pinned \texttt{amsart} editorial wrapper at 10pt on A4, which restates the theorem-like environments the retained text needs and adds the article's own macro definitions that the retained text needs (\texttt{\textbackslash Cc}, \texttt{\textbackslash epsilon}, \texttt{\textbackslash inv}), with extra wrapper packages (bm, bbm, dsfont, xspace, cancel, mathtools, enumitem). The delivered numbering is the unit's own. Assets: no retained span contains a figure environment, a graphics-inclusion command, a PDF-page inclusion, a table or a TikZ picture; no image file is copied into this package, the package declares no asset dependency and no image file is redistributed with it. Licence: version arXiv:2505.04296v2 of this article is distributed under the Creative Commons Attribution 4.0 International licence, as recorded in the arXiv licence metadata for this exact version and shown on the abstract page https://arxiv.org/abs/2505.04296v2. A full rights notice is delivered at licenses/arxiv-2505.04296-CC-BY-4.0.txt. Characters: every retained excerpt is pure ASCII, so the delivered engine, XeLaTeX with its default Latin Modern text font, reports no missing character.
\begin{prop}\label{coroll1}
Let $f(t) = t^m + a_{m-1}t^{m-1} + \dots + a_0$ and $g(t) = t^n + b_{n-1}t^{n-1} + \dots + b_0$ be polynomials in the ring $\Cc[t]$ with multisets of roots $[\lambda_j]$ and $[\mu_k]$, respectively. Then the characteristic polynomial
$$\chi(F(a_{n-1}, \dots, a_0)\boxplus F(b_{n-1}, \dots, b_0); t) = \det(t\cdot I - F(a_{n-1}, \dots, a_0)\boxplus F(b_{n-1}, \dots, b_0))$$
in the variable $t$ of the Kronecker sum of their Frobenius companion matrices has as its roots all possible $mn$ sums of the form $\lambda_j + \mu_k$.
\end{prop}

\begin{equation}\label{p_n}
p_n(z; x,y) = \prod\limits_{w\in \inv(x)\ast \inv(y)} (z - w),
\end{equation}

\begin{prop}\label{prop10}
For each positive integer number $n$, the following statements hold:

\begin{enumerate}[(i)]

\item The polynomial $p_n(z; x, y)$ from (\ref{p_n}) is symmetric in the variables $x$, $y$, and $z$.

\item For all complex $x$, $y$, and $z$, we have the equality
$$p_n(z^n; x, y) = \chi(F(\underset{n}{\underbrace{0,\dots,0,(-1)^{n+1}x}})\boxplus F(\underset{n}{\underbrace{0,\dots,0,(-1)^{n+1}y}}); z).$$

\item The polynomial $p_n(z; x, y)$ is a homogeneous polynomial in the variables $x$, $y$, and $z$ of degree $n$ with integer coefficients.
\end{enumerate}

\end{prop}

\begin{proof}
Expanding each factor in expression (\ref{p_n}) into $n$ linear terms, we obtain
\begin{equation}\label{eq1}
p_n(z; x, y) = \prod\limits_{r, s = 1}(\sqrt[n]{z} + \epsilon^r\sqrt[n]{x} + \epsilon^s\sqrt[n]{y}),
\end{equation}
where $\varepsilon$ is an (arbitrary) primitive $n$-th root of unity.

\begin{enumerate}[(i)]

\item Recall that we fix the root branch such that $\sqrt[n]{1} = 1$. From the equality (\ref{eq1}), it follows that the expression $p_n(z; x, y)$ is symmetric. Indeed,
\begin{align*}
p_n(x; z, y) &= \prod\limits_{r,s = 1}^n (\sqrt[n]{x} - \epsilon^r\sqrt[n]{(-1)^nz} - \epsilon^s\sqrt[n]{(-1)^ny}) \\
&= \prod\limits_{r,s = 1}^n (\sqrt[n]{x} + \epsilon^r\sqrt[n]{z} + \epsilon^s\sqrt[n]{y}) \\
&= \left(\epsilon \dots \epsilon^{n-1} \right)^n\prod\limits_{r, s =1}^n \epsilon^{-r}(\sqrt[n]{z} + \epsilon^{-r}\sqrt[n]{z} + \epsilon^{s+r}\sqrt[n]{y})\\
&= \prod\limits_{r,s = 1}^n (\sqrt[n]{z} + \epsilon^r\sqrt[n]{x} + \epsilon^s\sqrt[n]{y}) \\
&= p_n(z; x, y),
\end{align*}
where the second equality follows from the fact that $(-\epsilon^s\sqrt[n]{(-1)^nx})^n = x$, i.e., for any $n$, the expression $-\epsilon^s\sqrt[n]{(-1)^nx}$ runs over all roots of the polynomial $t^n - x$ as $s = 1, \dots, n$.

\item Note that from equation (\ref{eq1}), it also follows that $p_n(z; x, y)$ has as its roots all sums of the form $x_j + y_k$, where $\{x_j\}$ and $\{y_k\}$ are the sets of roots of the polynomials $z^n + (-1)^{n+1}x$ and $z^n + (-1)^{n+1}y$. Then from Proposition \ref{coroll1}, it follows that
$$p_n(z^n; x, y) =  \chi(F(\underset{n}{\underbrace{0,\dots,0,(-1)^{n+1}x}})\boxplus F(\underset{n}{\underbrace{0,\dots,0,(-1)^{n+1}y}}); z).$$

\item The statement follows immediately from point (ii) and the fact that the eigenvalues of the matrix $A^n$ are the $n$-th powers of the eigenvalues of the matrix $A$.
\end{enumerate}
\end{proof}

\end{document}
